% ACRONIMI SEMPLICI
\newacronym{bi}{BI}{Business Intelligence}
\newacronym{samm}{SAMM}{Software Assurance Maturity Model}


% TERMINI CON ACRONIMO + GLOSSARIO
% GDPR
\newacronym[description={\glslink{gdprg}{General Data Protection Regulation}}]
    {gdpr}{GDPR}{General Data Protection Regulation}

\newglossaryentry{gdprg} {
    name=\glslink{gdpr}{GDPR},
    text=General Data Protection Regulation,
    sort=gdpr,
    description={Regolamento generale sulla protezione dei dati personali emanato dall'Unione Europea (Regolamento UE 2016/679). Stabilisce le norme relative alla raccolta, al trattamento, alla conservazione e alla cancellazione dei dati personali, imponendo alle organizzazioni precisi obblighi in materia di trasparenza, consenso e sicurezza dei dati.}
}

% PII
\newacronym[description={\glslink{piig}{Personally Identifiable Information}}]
    {pii}{PII}{Personally Identifiable Information}

\newglossaryentry{piig} {
    name=\glslink{pii}{PII},
    text=Personally Identifiable Information,
    sort=pii,
    description={Qualsiasi dato che consenta di identificare direttamente o indirettamente una persona fisica (es. nome, cognome, indirizzo e-mail, numero di telefono, codice fiscale). Nel contesto del Data Catalog, i dati PII richiedono una gestione particolare in conformità alle normative vigenti.}
}

% RBAC
\newacronym[description={\glslink{rbacg}{Role-Based Access Control}}]
    {rbac}{RBAC}{Role-Based Access Control}

\newglossaryentry{rbacg} {
    name=\glslink{rbac}{RBAC},
    text=Role-Based Access Control,
    sort=rbac,
    description={Modello di controllo degli accessi che assegna i permessi agli utenti in base al ruolo che ricoprono nell'organizzazione, anziché individualmente. In piattaforme come OpenMetadata, il sistema RBAC determina quali operazioni possono essere eseguite su quali risorse.}
}


% VOCI DI GLOSSARIO
\newglossaryentry{anonimizzazione}{
    name={Anonimizzazione},
    description={Processo irreversibile mediante il quale i dati personali vengono alterati in modo tale da non consentire l'identificazione, diretta o indiretta, della persona fisica a cui si riferiscono. Una volta anonimizzati, i dati non sono più soggetti alle disposizioni del GDPR.}
}

\newglossaryentry{pseudonimizzazione}{
    name={Pseudonimizzazione},
    description={Processo mediante il quale i dati personali vengono trattati in modo tale da non poter essere attribuiti a una persona fisica specifica senza l'utilizzo di informazioni aggiuntive. A differenza dell'anonimizzazione, la pseudonimizzazione è un processo reversibile e i dati rimangono soggetti al GDPR.}
}

\newglossaryentry{dataquality}{
    name={Data Quality},
    description={Misura che valuta quanto un insieme di dati sia "sano" e adatto allo scopo per cui deve essere utilizzato dall'azienda. Applicare la Data Quality significa adottare un sistema che implementa controlli automatici sui dati per individuare anomalie, dati incompleti o informazioni incoerenti.}
}

\newglossaryentry{datagovernance}{
    name={Data Governance},
    description={Insieme di policy, ruoli, regole e processi aziendali che garantiscono che i dati siano sicuri, affidabili, accessibili e conformi alle normative. Rappresenta il "regolamento aziendale" che gestisce il ciclo di vita del dato.}
}

\newglossaryentry{dataingestion}{
    name={Data Ingestion},
    description={Processo di estrazione e importazione automatica dei metadati dalle fonti dati collegate al catalogo. Include la raccolta di informazioni strutturali (tabelle, colonne) ed eventuale profilazione e classificazione dei dati.}
}

\newglossaryentry{datalineage}{
    name={Data Lineage},
    description={Rappresenta la tracciabilità del dato lungo il suo intero ciclo di vita: indica l'origine del dato, le trasformazioni che ha subito e i sistemi attraverso cui è transitato prima di giungere alla destinazione finale, aumentando la fiducia nell'affidabilità delle informazioni.}
}

\newglossaryentry{dataobservability}{
    name={Data Observability},
    description={Capacità di conoscere in tempo reale lo "stato di salute" dei dati. Applica i concetti di monitoraggio del software ai flussi di dati, permettendo di identificare, diagnosticare e risolvere i problemi prima che impattino gli utenti aziendali.}
}

\newglossaryentry{datadiscovery}{
    name={Data Discovery},
    description={Processo di ricerca e centralizzazione volto a scoprire quali informazioni sono presenti all'interno dei vari sistemi dati aziendali. In una piattaforma di Data Catalog, agisce come un motore di ricerca intuitivo per esplorare i dataset disponibili.}
}
